% Ciclo de vida de los contenedores
\subsection{Ciclo de vida de los contenedores}\label{sec:6_4_ciclo_vida}

\subsubsection{Reinicio automático (-{}-restart always)}
Los once contenedores del Backend quedan con la política \texttt{-{}-restart always}: Podman los vuelve a iniciar si el proceso termina (por ejemplo, por falta de memoria) y el servicio \texttt{podman-restart} los levanta al encender la VPS. Se comprueba para todos a la vez:
\begin{lstlisting}
# podman inspect --format '{{.Name}} {{.HostConfig.RestartPolicy.Name}}' $(podman ps -aq)
\end{lstlisting}
Cada línea debe terminar en \texttt{always}.
\captura{back-08-restart-policy}{Política de reinicio de los contenedores del Backend}

\subsubsection{Redespliegue de un servicio (up-container y delete-container)}
Para actualizar un microservicio se sustituye su JAR y se vuelve a construir solo ese servicio; los demás siguen funcionando:
\begin{enumerate}
  \item Compilar el servicio en el equipo de desarrollo y copiar el JAR a \texttt{artifact/} con el mismo nombre (secciones \ref{sec:4_2_compilacion_local} y \ref{sec:4_4_transferencia_vps}).
  \item En el VPS Backend:
\begin{lstlisting}
# cd /apps/aplicaciones/defensoria/back
# bash podman-compose.sh up-container quejas-service
# podman logs -f --tail 50 quejas-service        # Ctrl+C para salir
\end{lstlisting}
\end{enumerate}
\texttt{up-container} elimina el contenedor, reconstruye la imagen con el JAR nuevo y crea el contenedor otra vez. Durante ese tiempo (menos de un minuto) solo ese servicio deja de responder.

\begin{table}[H]
\centering
\caption{Operaciones del ciclo de vida}
\small
\begin{tabular}{p{4.6cm}p{9.8cm}}
\toprule
\textbf{Situación} & \textbf{Comando} \\
\midrule
Cambió el JAR & \texttt{bash podman-compose.sh up-container <servicio>} \\
Cambió solo el config-file & \texttt{podman restart <servicio>} \\
Detener un servicio & \texttt{bash podman-compose.sh delete-container <servicio>} \\
Detener todo el Backend & \texttt{bash podman-compose.sh delete} y \texttt{bash podman-compose-modelo.sh delete} \\
Actualizar el modelo de IA & \texttt{bash podman-compose-modelo.sh up} \\
Liberar espacio de imágenes viejas & \texttt{podman image prune} \\
\bottomrule
\end{tabular}
\end{table}
\captura{back-09-up-container}{Redespliegue de un microservicio con up-container}
